\section{Численные методы решения}

\subsection{Метод Эйлера}

Метод Эйлера является простейшим численным методом решения обыкновенных дифференциальных уравнений. Для системы первого порядка \(\dot{y} = f(t, y)\) итерационная схема имеет вид:

\begin{equation}
    y_{n+1} = y_n + h \cdot f(t_n, y_n)
    \label{eq:euler}
\end{equation}

где \(h\) — шаг интегрирования. Для системы второго порядка (уравнения движения) необходимо привести уравнения к системе первого порядка.

\subsection{Метод Рунге-Кутты 4-го порядка}

Метод Рунге-Кутты 4-го порядка (RK4) обеспечивает значительно более высокую точность по сравнению с методом Эйлера \cite{kalitkin2021}. Итерационная схема выглядит следующим образом:

\begin{equation}
\begin{aligned}
    y_{n+1} &= y_n + \frac{h}{6}(k_1 + 2k_2 + 2k_3 + k_4) \\
    k_1 &= f(t_n, y_n) \\
    k_2 &= f(t_n + \frac{h}{2}, y_n + \frac{h}{2}k_1) \\
    k_3 &= f(t_n + \frac{h}{2}, y_n + \frac{h}{2}k_2) \\
    k_4 &= f(t_n + h, y_n + h k_3)
\end{aligned}
\label{eq:rk4}
\end{equation}

\subsection{Программная реализация}

Ниже представлена реализация численных методов на языке \texttt{Python} с использованием библиотек \texttt{NumPy} и \texttt{Matplotlib}. Полный код программы приведен в листинге~\ref{code:python}.

\inputminted[frame=lines, framesep=5mm, baselinestretch=1.2, fontsize=\small, linenos]{python}{code/motion_simulation.py}

Ключевые элементы программы:

\begin{itemize}
    \item Функция \mintinline{python}{system(t, state, with_drag)} вычисляет производные состояния системы;
    \item Функция \mintinline{python}{euler_method(...)} реализует метод Эйлера \eqref{eq:euler};
    \item Функция \mintinline{python}{rk4_method(...)} реализует метод Рунге-Кутты 4-го порядка \eqref{eq:rk4};
    \item Переменная \texttt{with\_drag} управляет учетом сопротивления воздуха.
\end{itemize}